\section{Prompts}
\label{app:prompts}
In this section, we provide the specific prompts used in the DynaDebate framework. To provide a concrete demonstration of the instruction design, we use the MATH500 dataset as an illustrative example throughout this section.

\subsection{Path Generation and Allocation}
\label{app:prompt_path}

\begin{tcolorbox}[
    breakable,
    colback=white,
    colframe=black,
    boxrule=0.8pt,
    left=6pt,
    right=6pt,
    top=6pt,
    bottom=6pt
]
\ttfamily\small
System Instruction:\\
You are participating in a collaborative problem-solving session for a challenging math competition problem.

\vspace{0.5em}
--- STRATEGIC BRAINSTORMING PHASE ---\\
Your task is to identify genuinely independent and feasible mathematical approaches that could solve the problem.
A ``path'' here refers to a distinct mathematical strategy, such as a specific logical approach, theorem application,
or conceptual framework (e.g., algebraic vs.\ geometric, recursion vs.\ combinatorics).

\vspace{0.5em}
Problem: \lbrack Problem Description\rbrack

\vspace{0.5em}
Your responsibilities:

\vspace{0.25em}
PART 1: METHOD EXPLORATION\\
1. List up to 3 viable candidate strategies that could be used to solve the question.\\
2. List ONLY genuinely distinct strategies.
\begin{itemize}
    \item If only 1 or 2 viable strategies exist, list only those.
    \item Do NOT invent additional strategies merely to increase the count.
    \item Strategies that differ only in notation or minor algebraic manipulation but rely on the same core theorem
    do NOT count as independent.
\end{itemize}

\vspace{0.25em}
PART 2: RIGOROUS FEASIBILITY ASSESSMENT\\
For each strategy, specify:
\begin{itemize}
    \item Core idea (key theorem, formula, or logical technique)
    \item Expected reliability (High / Medium / Low)
\end{itemize}

\vspace{0.25em}
PART 3: COLLABORATION STRATEGY RECOMMENDATION\\
Output strictly in the following format:

\begin{verbatim}
"METHOD_1:"
- Core idea: ...
- Critical step: ...

"METHOD_2:"
...
\end{verbatim}
\end{tcolorbox}

\subsection{Process-Centric Debate}
\label{app:prompt_debate}

\begin{tcolorbox}[
    breakable,
    colback=white,
    colframe=black,
    boxrule=0.8pt,
    left=6pt,
    right=6pt,
    top=6pt,
    bottom=6pt
]
\ttfamily\small
System Instruction:\\
Please solve the following high school competition math problem. Provide a step-by-step reasoning for your solution.
Use LaTeX for all mathematical expressions.

\vspace{0.5em}
Problem: \lbrack Problem Description\rbrack

\vspace{0.5em}
--- COMPARISON PHASE (Round N) ---

\vspace{0.25em}
Your Original Method: \lbrack Method Name\rbrack: \lbrack Method Description\rbrack\\
You are Agent \lbrack ID\rbrack.

\vspace{0.25em}
Previous Round Results from All Agents:\\
Agent 1 Result: [Response 1]\\
...

\vspace{0.25em}
[Optional: Code Verification Feedback Inserted Here]

\vspace{0.5em}
--- YOUR TASK ---

\vspace{0.25em}
Objective: Compare results, consider verification feedback, and refine your solution.

\vspace{0.25em}
Analysis Required:\\
1. Critically review the reasoning from other agents. Do you see any logical flaws or calculation mistakes?\\
2. Consider the computational verification results. Re-evaluate your own previous reasoning in light of their approaches.\\
3. Re-evaluate your own previous reasoning in light of their approaches.

\vspace{0.5em}
Given a mathematics problem, determine the answer through comparison and verification feedback.
Your final answer should be in the form \verb|\boxed{answer}|, at the end of your response.
\end{tcolorbox}
\label{fig:prompt_debate}

\subsection{Verification Mechanism}
\label{app:prompt_verification}

\begin{tcolorbox}[
    breakable,
    colback=white,
    colframe=black,
    boxrule=0.8pt,
    left=6pt,
    right=6pt,
    top=6pt,
    bottom=6pt
]
\ttfamily\small
System Instruction:\\
You are in a multi-agent collaborative setting for solving a math problem, and your task is to generate code to verify the answers.

\vspace{0.5em}
--- PROBLEM --- \\
\lbrack Problem Description\rbrack

\vspace{0.5em}
--- VERIFICATION TRIGGER --- \\
Reason for verification: \lbrack Trigger Reason\rbrack

\vspace{0.5em}
--- RESPONSES TO VERIFY --- \\
Agent 1 Response: ...\\
Agent 2 Response: ...

\vspace{0.5em}
--- YOUR TASK ---

\vspace{0.25em}
Objective: Use computational verification to analyze the responses and provide feedback for the next round of debate.

\vspace{0.25em}
Verification Strategy:\\
1. Extract Key Claims: Identify the main mathematical claims and final answers.\\
2. Design Verification Code: Write Python code to solve the current problem.\\
3. Execute and Compare: Run your code and compare results with agent responses.\\
4. Identify Issues: Determine which responses (if any) contain errors.\\
5. Provide Constructive Feedback: Give specific guidance for the next round.

\vspace{0.5em}
CRITICAL REQUIREMENTS FOR CODE:
\begin{itemize}
    \item Store your final result in a variable named `ans'
    \item Write clean, readable Python code
    \item Use appropriate mathematical libraries (math, numpy, sympy, etc.)
\end{itemize}

\vspace{0.25em}
Your response must include:\\
1. Analysis Summary: Brief explanation of discrepancies or concerns found.\\
2. Verification Code: Python code in python blocks.\\
3. Agent-Specific Feedback: Specific feedback for each agent.\\
4. Recommendations: Concrete suggestions for improving solutions.

\vspace{0.5em}
Format your code properly. This feedback will be provided to all agents in the next round.
\end{tcolorbox}
\label{fig:prompt_verification}
